Each step names the result that licenses it; the derivation is those results.

The Principle says what is true; practice needs to know what to do. The theorems distil into a six-step procedure, and every step names the theorem, proved above, that licenses it. Table \ref{tab:algorithm} is the one-page version: the theorems in operational form. That the steps form a decision procedure with no gaps is itself a theorem, Corollary \claimref{transport.no-fourth-rule}.




\begin{enumerate}
\item[\textbf{Step 1.}] \textbf{Locate the failure.} Identify the element at which invertibility fails, or, if the construction is a family, the locus at which it fails. Nothing later depends on how this is done; it is a computation in the case at hand. What the algebra guarantees is that the failure cannot be argued away: a null vector has no inverse in \emph{any} unital algebra containing the base (Lemma \claimref{transport.noinverse}), and neither has an idempotent.

\item[\textbf{Step 2.}] \textbf{Diagnose the type.} Ask, in order:
  \begin{itemize}
  \item Is the failing element a vector $v$ with $Q(v) = 0$? Then the obstruction is \textbf{type N}. These exist exactly in mixed signature, $pq \neq 0$ (Proposition \claimref{transport.existence-locus-type-n}).
  \item Is it a \emph{central} idempotent? Then \textbf{type C}. These exist exactly when $p - q \equiv 1 \pmod 4$, and are then exactly the two projectors $\tfrac12(1 \pm \omega)$ (Theorem \claimref{transport.existence-classification-type-c}); equivalently, exactly when the algebra is not simple (Theorem \claimref{transport.classification-real-clifford-algebras}).
  \item Otherwise \textbf{type P}: a non-central idempotent, or an element annihilated by one.
  \end{itemize}
  The trichotomy is exhaustive: every zero divisor of a nondegenerate real Clifford algebra is annihilated by a nontrivial idempotent, possibly central in the split signatures, and then exactly when the element lies in one ideal (Theorem \claimref{transport.exhaustiveness-every-zero-divisor-annihilated-idempotent}). Two cautions. The nilpotency locus is strictly larger than the type-N locus, and nilpotence alone does not diagnose N (Example \ref{ex:cl20}); and a single algebra may carry all three types at once (Example \ref{ex:cl10}).

\item[\textbf{Step 3.}] \textbf{Apply the rule.}
  \begin{itemize}
  \item \textbf{Type N, Rule 1.} Adjoin one hyperbolic plane: $\Cl(p,q) \to \Cl(p{+}1,q{+}1)$. Do not expect the element to be repaired: no equivariant assignment of a partner to $v$ exists, in this ascent or any other (Theorem \claimref{transport.no-go-level-elements}), and the obstruction is a cohomology class that every enlargement pulls back unchanged (Theorem \claimref{transport.cohomological-form-no-go}). What is repaired is the space of points (Step 4). The increment is forced: $(+1,+1)$, never $(+2,0)$ or $(0,+2)$, among null-lift completions (Theorem \claimref{transport.minimality}), and no algebra with a smaller or differently balanced quadric can host the resolution however the base is embedded in it (Theorem \claimref{transport.representation-independent-minimality-completion}).
  \item \textbf{Type C, Rule 2.} Adjoin one generator $\nu$ orthogonal to $\mathbb{V}$, of either nonzero square. The grade involution becomes inner, the two ideals fuse, and the completion is $\Mat_2(\Cl^0(p,q))$, simple (Theorem \claimref{transport.rule-2-ideal-fusion-sharp-form}). The null-generator ascent is not an option: it produces a square-zero ideal and never fuses (Theorem \claimref{transport.one-generator-minimal-minimal-completions-classified}(2)). The sign of $\nu^2$ is free and the algebra forgets it (Remark \ref{rem:forgottenbit}); if it matters for the application, it must be fixed by structure outside the algebra.
  \item \textbf{Type P.} Do nothing. The identity assignment $\sigma(e) = e$ is a total, fully equivariant von Neumann repair, at signature increment $(0,0)$ (Proposition \claimref{transport.type-p-self-repairing}).
  \end{itemize}

\item[\textbf{Step 4.}] \textbf{Compute in the completion.} For type N this is the substantive step, and it is where the classical formulas apply. Represent points, not elements: lift $x \in \mathbb{V}$ to the null ray of
  \[
  X(x) \;=\; \nO + x + \tfrac12 Q(x)\,\nInf ,
  \]
  which satisfies the incidence condition $Q(x) = 2st$ in the homogeneous coordinates $(s,x,t)$ of Proposition \claimref{transport.swap-rescale-normal-form}. The reciprocal is then the linear involution
  \[
  \iota(s,\,x,\,t) \;=\; (2t,\ x,\ \tfrac12 s),
  \]
  total on the quadric and never failing (Theorem \claimref{transport.totality}). All nonlinearity and all division have been deferred to a single read-out. This is the linearise-then-normalise pattern of homogeneous coordinates and M\"obius matrices, encountered here as forced rather than chosen (Remark \ref{rem:classical}).

\item[\textbf{Step 5.}] \textbf{Descend, and state what descent costs.} Return the answer to $\mathbb{V}$ with the read-out $\dsc[\,s\nO + x + t\nInf\,] = x/s$. On the invertible locus the composite gives exactly the classical answer (Lemma \claimref{transport.descent-read-out-map}(2)); where the first slot vanishes the answer is an ideal point and corresponds to no point of $\mathbb{V}$ (Lemma \claimref{transport.descent-read-out-map}(3)). \emph{The refusal of the read-out is information, not failure}: it says the answer exists and lies outside the original space. At the level of elements the corresponding statement is stronger and worth knowing before relying on any descent: the pinching onto the base is a conditional expectation whose kernel is exactly the adjoined ladder sector (Proposition \claimref{transport.pinching-kernel}), and there is no unital algebra homomorphism from either completion back to its base at all (Theorem \claimref{transport.door-one-way}). Descent is invertible on the transported elements and impossible on the elements the ascent adjoined.

\item[\textbf{Step 6.}] \textbf{Iterate, or stop.} Ask whether the completion carries an obstruction of its own. It does: no transport terminates in a division algebra (Proposition \claimref{transport.no-transport-terminates}). But the bookkeeping is finite in the following sense: Rule 1 fixes the mod-8 residue, so the taxonomy of the completion is that of the base and iterating changes nothing structural (Theorem \claimref{transport.rule-1-morita-trivial}), while Rule 2 moves the residue off the split classes and can therefore occur at most once in any transport history (Proposition \claimref{transport.fusion-occurs-most-once}). A transport history takes at most one step on the clock, and then remains fixed.
\end{enumerate}
